\section{Contribution to the Dissertation}
\label{sec:receipts-contribution}

\subsection{Contribution to Prediction vs.~Coordination}
\label{sec:receipts-contribution-prediction}

The central epistemological tension in the doctoral thesis is between prediction and
coordination. A predictive system makes claims about what will happen based on learned
patterns. A coordination system makes claims about what has happened based on observed
event evidence, and uses those claims to coordinate action. The 2016 language-model
lesson established that prediction alone---local text imitation---does not conserve
consequence. The receipts corpus establishes the complementary principle: that
coordination claims, to be board-admissible, must be grounded in receipted event evidence.

The five M\&A claim receipts are coordination artifacts, not prediction artifacts. Each
receipt documents a conformance measurement computed over an observed event log---a
record of what a process actually did, not a model of what it is predicted to do.
The fitness of $0.982$ in \texttt{rec\_ebitda\_rework\_001.json} is not a prediction
of EBITDA improvement; it is a measurement of how closely the observed manufacturing
process conforms to a standard process model, from which the \$1.25M EBITDA impact
is derived. This distinction between prediction and measurement is the epistemological
core of the receipts corpus's contribution to the dissertation.

The receipts corpus demonstrates that coordination at board level---the coordination
of capital allocation decisions in an M\&A transaction---requires an evidence chain
that cannot be manufactured from prediction alone. It requires receipted, hash-bound,
signed measurements derived from observed event logs. This is the empirical refutation
of the old-regime assumption identified in Section~\ref{sec:receipts-question-old-regime}.

\subsection{Contribution to Process Evidence}
\label{sec:receipts-contribution-evidence}

The receipts corpus is the primary contribution of this research program to the theory
and practice of process evidence. Process evidence, as defined in the Van der Aalst
Constitution (Process Mining Chicago TDD), is event log evidence that can be mined into
a conforming object-centric process. The receipts corpus operationalizes this definition
by requiring that every board-admissible claim be traceable to a specific event log
identified by a BLAKE3 hash.

The corpus demonstrates three properties of process evidence that the dissertation
argues are necessary for board admissibility. First, \emph{content-addressability}:
the BLAKE3 hash ensures that the specific log from which a claim was derived can be
identified and re-accessed. Second, \emph{measurement reproducibility}: the
\texttt{wasm4pm\_query\_result} payload in each receipt documents the specific algorithm,
parameters, and numerical result of the conformance computation, allowing independent
replication. The \texttt{rec\_wc\_ar\_002.json} receipt explicitly carries a
replication verdict of \texttt{PASSED}, demonstrating that this property has been
verified for at least one of the five claims. Third, \emph{tamper-evidence}: the
Ed25519 signature over the JCS canonical form of the entire JSON document ensures
that no field can be modified after signing without detection.

Together these three properties define what the dissertation terms a \emph{receipted
fact}---a claim that has crossed the manufacture-to-claim boundary with its evidence
chain intact and its integrity sealed. The receipts corpus is the first instantiation
of this concept in the process-intelligence research program.

\subsection{Contribution to Manufacturing Architecture}
\label{sec:receipts-contribution-manufacturing}

The receipts corpus makes a structural contribution to the manufacturing architecture
of the \texttt{wasm4pm} execution platform. The module generation receipts establish
the principle that the manufacturing pipeline itself must be receipted, not just its
outputs. A conformance authority module that produces correct results but whose
manufacturing provenance is unknown cannot be trusted to produce correct results
in a future context. The generation receipt for version \texttt{v30.1.2}---Blake3
hash \texttt{75508e6e9e7acccb679e5c8be35ddd2adc9c4732f1fc331047d755703c9107c3},
$8/8$ tests passing, zero compilation errors---establishes the manufacturing provenance
of the module that underwrites the five M\&A claim receipts.

This recursive receipting---receipts for the receipting infrastructure---is the
architectural contribution that distinguishes the receipts corpus from conventional
audit trails. It implements a trust chain that extends from the board presentation
artifacts (Tera2 templates, \texttt{pptx-rs} rendering) backward through the
conformance authority module, through the \texttt{cargo check} and \texttt{cargo test}
results, to the Rust source code identified by its Blake3 hash. At every stage of
this chain, the evidence is bound to a specific, immutable artifact. No stage can
be silently substituted without breaking the hash chain.

\subsection{Contribution to Capital-Flow Infrastructure}
\label{sec:receipts-contribution-capital}

The doctoral lineage culminates in a capital-flow and settlement infrastructure that
depends on receipted process evidence as its legitimacy layer. The thesis argues that
autonomous capital allocation---whether in M\&A transactions, DAO governance, or
AI-mediated settlement---requires a receipted evidence chain of the kind demonstrated
by this corpus. A DAO voting on a capital allocation decision cannot rely on asserted
metrics; it requires receipted measurements that any member can independently verify.

The receipts corpus demonstrates the feasibility of this infrastructure at the
M\&A due diligence scale. The five board-admissible claim receipts, totalling
approximately \$3.07M in documented process-improvement value across EBITDA rework,
working capital velocity, SLA penalty exposure, compliance leakage, and residual
standardization, constitute a capital-flow claim set that is fully receipted, fully
hash-bound, and fully signed. This is the proof of concept for a capital-flow
infrastructure in which every value claim carries its own cryptographic derivation
chain.

The \texttt{RECEIPT\_REGISTRY.md} anticipates the extension of this infrastructure
through its \texttt{ADVERSARIAL} receipt type, which is reserved for claims that
have been stress-tested against hostile process variants. This anticipates a future
in which capital-flow claims must pass not only conformance gates but also adversarial
resilience gates, a requirement that becomes essential as the infrastructure moves
from M\&A due diligence into live settlement and DAO governance contexts.
